\subsection*{Problems}

\textbf{6.1.} Consider $\{0,1,2,\dots,n\}$ as a sample space with the discrete $\sigma$-algebra. Let $p$ be a real number between 0 and 1. Define a probability measure $P$ by
\[
P(\{k\}) = \binom{n}{k}p^k(1-p)^{n-k},
\]
for $k=0,1,2,\dots,n$. Evaluate the expectation of the following random variables:

\noindent (a) \quad $X(k)=k$.

\noindent (b) \quad $Y(k)=
\begin{cases}
1 & \text{if } k \text{ is even}, \\
-1 & \text{if } k \text{ is odd}.
\end{cases}$

\textbf{6.2.} In the hat problem described in Example 6.5.1, find the variance of the number of people who can get back his/her own hat. (Hint: Compute $E[1_{A_i}1_{A_j}]$.)

\textbf{6.3.} There are $n$ different types of coupons. In each iteration, a coupon of random type is drawn. Let $X_k$ denote the time you first get $k$ different types of coupons. Trivially, we have $X_1=1$. $X_2$ is the first time that you get a coupon different from the first one. For $k \ge 1$, find the expectation of $X_k$.

\textbf{6.4.} Let $X$ be a nonnegative random variable that may take $\infty$ as its value. Show that if $E[X]$ is finite, then $X$ is finite almost surely.

\textbf{6.5.} Let $X$ be an integrable random variable. For any $\epsilon$, show that there is a simple function $f$ such that $E[|X-f|]<\epsilon$.

\textbf{6.6.} Suppose $X$ is a random variable in $L^1(P)$. Show that
\[
\lim_{n\to\infty} \int X1_{|X|\ge n}\, dP = 0.
\]

\textbf{6.7.} Consider a sequence $X_1,X_2,X_3,\dots$ of nonnegative independent random variables, where $P(X_k>a)\ge \delta$ for all $k$. Here, $a$ is a positive constant and $\delta$ is another positive constant between 0 and 1. Show that $\sum_{k=1}^{\infty} X_k=\infty$ with probability 1. (Hint: Use the Borel--Cantelli lemma.)

\textbf{6.8.} Let $(\Omega,\mathcal{F},\mu)$ be a measure space and $X$ be a nonnegative measurable function defined on this measure space such that $\int_{\Omega} X\, d\mu =1$. Prove that the set function $Q$, defined by $Q(A)=\int_A X\, d\mu$ for $A\in\mathcal{F}$, is a probability measure.

\textbf{6.9.} The objective of this question is to prove that for any infinite array of nonnegative real numbers $a_{ij}$, for $i\ge 1$ and $j\ge 1$, the two different orders of summation give the same answer
\begin{equation}
\sum_{i=1}^{\infty}\sum_{j=1}^{\infty} a_{ij}
=
\sum_{j=1}^{\infty}\sum_{i=1}^{\infty} a_{ij}.
\tag{6.5}
\end{equation}

(The answer may equal $\infty$. In this case, both sides are equal to $\infty$.)

To formulate the problem using the measure theory, we take $\mathbb{N}=\{1,2,3,\dots\}$ as the sample space, equipped with the discrete $\sigma$-algebra $\mathcal{P}(\mathbb{N})$. Let $\mu$ denote the counting measure on $\mathbb{N}$, i.e., $\mu(A)$ is the size of $A$ for any $A\subseteq \mathbb{N}$.

\noindent (a) For each fixed $i\ge 1$, let $X_i$ be the function defined by $X_i(j)=a_{ij}$. Verify that
\[
\int X_i\, d\mu = \sum_{j=1}^{\infty} a_{ij}.
\]

\noindent (b) Let $Y_n=X_1+X_2+\cdots+X_n$. Prove that $Y_1,Y_2,Y_3,\dots$ is an increasing sequence of functions. What is the limit function $Y \triangleq \lim_n Y_n$?

\noindent (c) Apply the monotone convergence theorem to prove (6.5).

\textbf{6.10.} Consider the random series
\[
S = 1 \pm \frac{1}{2} \pm \frac{1}{4} \pm \frac{1}{8} \pm \frac{1}{16} \pm \frac{1}{32} \pm \cdots,
\]
where the signs are chosen independently by tossing a fair coin infinitely many times. Find the mean and the variance of $S$.
